\documentclass[11pt]{letter}
\usepackage[T1]{fontenc}
\usepackage{lmodern}
\usepackage[margin=2.5cm]{geometry}
\usepackage[hidelinks]{hyperref}
%% TODO(author): confirm name and address.
\signature{Subrahmanyam Tanala\\Department of Electrical and Electronics Engineering\\Anil Neerukonda Institute of Technology and Sciences (ANITS)\\Visakhapatnam, Andhra Pradesh, India\\\href{mailto:subrahmanyam.eee@anits.edu.in}{subrahmanyam.eee@anits.edu.in}}
\address{}
\begin{document}
\begin{letter}{The Editor-in-Chief\\Computers in Biology and Medicine}
\opening{Dear Editor,}

Please consider the enclosed manuscript, ``Leakage-controlled evaluation of stacking and recursive feature elimination for diabetes outcome prediction on the PIMA dataset'', for publication as a Research Article in \emph{Computers in Biology and Medicine}.

Stacking ensembles and wrapper feature selection are widely reported for diabetes prediction, including in this journal, often with very high accuracy on the PIMA Indians Diabetes dataset. Such results are hard to interpret when imputation or feature selection is fitted before validation and when differences between models are not tested statistically. Our study takes a methodological rather than an algorithmic approach:

\begin{itemize}
  \item every data-dependent step (imputation, feature engineering, scaling, RFECV and stacking) is re-fitted inside each validation split;
  \item eleven learners (the eight base learners of the stack, two reference learners and the stack itself) are compared with and without a common RFECV subset using 5$\times$10-fold repeated cross-validation, corrected confidence intervals and multiplicity-adjusted tests, and confirmed on an untouched hold-out set with DeLong tests and bootstrap intervals;
  \item the evaluation goes beyond discrimination to calibration, sensitivity-targeted operating thresholds, decision curves, a missing-data sensitivity analysis and SHAP explanations.
\end{itemize}

The main finding is deliberately reported without overstatement. RFECV produces compact, clinically plausible feature sets, and the stacking ensemble gives the highest observed hold-out sensitivity and MCC. However, it does not discriminate better than regularised logistic regression once validation is leakage-free, and its class-weighted probabilities are systematically too high, a problem that accuracy-based reporting would hide. The PIMA predictors include 2-h OGTT values, so we present the work as a methodological benchmark of internal validity, not as a deployable screening tool. We believe these results, together with fully reproducible code, will be useful to readers developing and reviewing clinical prediction models.

This manuscript is original, has not been published previously and is not under consideration elsewhere. The author has no competing interests to declare. All code and results are publicly available.

\closing{Yours sincerely,}
\end{letter}
\end{document}
